\documentclass[11pt]{article}
\usepackage[margin=0.95in]{geometry}
\usepackage{amsmath,amssymb,booktabs,array}
\usepackage{xcolor}

\newcommand{\M}{\mathcal M}
\newcommand{\T}{\mathsf{T}}

\newenvironment{item12}[1]{%
  \par\medskip\noindent\rule{\textwidth}{0.7pt}\par\nobreak\smallskip
  \noindent\textbf{\large #1}\par\nobreak\smallskip
}{\par\medskip}

\newcommand{\ask}[1]{\noindent\emph{Asked.}\ #1\par\smallskip}
\newcommand{\resp}{\noindent\textbf{Response.}\ }

\setlength{\parindent}{0pt}
\setlength{\parskip}{4pt}

\title{Responses to the twelve items}
\author{}
\date{27 August 2026}

\begin{document}
\maketitle

Answered in the order given. Items 6 and 11 contain corrections rather than confirmations, and
both are flagged at the point they arise. All numbers below are from completed runs; nothing in
this document is pending.

% =================================================================
\begin{item12}{1. Super-resolution distributional comparison: sample count, seeds, setup, and
the $96\times96$ results}
\ask{Confirm the evaluation sample count and number of seeds, confirm whether the preprocessing
and train/test setup match the dimensionality experiment, and add the $96\times96$ results.}

\resp Every arm is evaluated on the full STL-10 test split, $1{,}024$ images, with three
training seeds per arm ($0,1,2$); the numbers below are means over those three seeds, with the
seed standard deviation given for PSNR.

The setup is not merely matched to the dimensionality experiment, it is the same code path: both
call one loader that returns $5{,}000$ training and $1{,}024$ test images, resized to the target
resolution, centre-cropped, and normalised to $[-1,1]$. The dimensionality measurement is taken
inside the same loop, before any training, so no discrepancy is possible.

The $96\times96$ arm has completed. All three resolutions:

\begin{center}
\small
\begin{tabular}{@{}llrrrr@{}}
\toprule
res & arm & PSNR $\uparrow$ & energy $\downarrow$ & $W_1^{\mathrm{hf}}$ $\downarrow$ &
restored detail\\
\midrule
$32$ & identity            & $19.7421$ & $6.7433$ & $154.02$ & $0.000$\\
     & regression          & $20.5045$ & $0.1200$ & $49.37$  & $0.680$\\
     & $\M_1$ marginal     & $21.4927$ & $0.7335$ & $92.82$  & $0.397$\\
     & $\M_2$ source-cond. & $21.5575$ & $1.0168$ & $104.09$ & $0.324$\\
     & \emph{null floor}   &           & $-0.0029$ & $7.04$  & \\
\midrule
$48$ & identity            & $20.7718$ & $8.9424$ & $270.13$ & $0.000$\\
     & regression          & $21.7829$ & $0.0799$ & $67.69$  & $0.750$\\
     & $\M_1$ marginal     & $22.9923$ & $0.8259$ & $154.89$ & $0.427$\\
     & $\M_2$ source-cond. & $23.0912$ & $0.9595$ & $164.48$ & $0.391$\\
     & \emph{null floor}   &           & $0.0012$ & $12.32$  & \\
\midrule
$96$ & identity            & $21.0806$ & $17.1672$ & $998.74$ & $0.000$\\
     & regression          & $21.7075$ & $0.1023$  & $215.45$ & $0.784$\\
     & $\M_1$ marginal     & $23.1984$ & $1.8884$  & $623.56$ & $0.376$\\
     & $\M_2$ source-cond. & $23.1984$ & $2.2222$  & $659.37$ & $0.340$\\
     & \emph{null floor}   &           & $0.0010$  & $17.63$  & \\
\bottomrule
\end{tabular}
\end{center}

``Energy'' is the energy distance between the distribution of null-space residuals --- the
detail the low-resolution source does not determine --- of the outputs and of the targets.
$W_1^{\mathrm{hf}}$ is the one-dimensional Wasserstein distance between the distributions of
high-frequency energy. The null floor for both is obtained by splitting the target set in half
and comparing the halves, so each column has its own zero.

At every resolution $\M_2$ attains the higher or equal PSNR and the worse value on both
distributional columns, in each of the three seeds separately. At $96\times96$ the two models are
indistinguishable on PSNR to four decimals ($23.1984$ against $23.1984$, seed spreads $0.018$ and
$0.052$) while $\M_2$ remains $18\%$ worse on the energy distance. The comparison is drawn
between $\M_1$ and $\M_2$, which are matched in architecture, capacity and training budget; the
regression arm is included for reference only, and its high restored-detail fraction alongside
its low PSNR shows that the fraction counts incorrect detail as well as correct detail.
\end{item12}

% =================================================================
\begin{item12}{2. Parameterisation of the transport strengths, and what ``retaining'' and
``transporting'' mean}
\ask{Specify how the seven transport strengths are parameterised in the Stable Diffusion v1.5
pipeline, and define operationally what is meant by retaining the source versus transporting it.}

\resp The strength is the image-to-image pipeline's own \texttt{strength} argument,
$s\in\{0.1,0.2,0.3,0.4,0.5,0.6,0.8\}$. Operationally: the degraded image is encoded to a latent,
noised to the timestep a fraction $s$ of the way along the schedule, and denoised from there to
the end. With a $30$-step schedule, $\lceil 30 s\rceil$ steps are actually executed. $s$ is
therefore the fraction of the generative trajectory that is run: at $s\to0$ the pipeline returns
its input, at $s=1$ it discards the input and samples from noise. Guidance scale is $7.5$ and the
prompt is ``a photo of a \{class\}'' throughout.

\emph{Retaining the source} means returning the degraded image itself, unmodified. It is the
baseline row, and its gain is zero by definition.
\emph{Transporting} means returning the pipeline output at strength $s$.

Both are scored the same way: the CLIP cosine between the embedding of the returned image and the
embedding of the clean original. The reported gain is the transported score minus the retained
score, so a negative gain means the pipeline made the image less like its own original than
leaving it alone would have.

Degradation is downsample-then-upsample by a factor $f\in\{2,4,8,16\}$ on $64$ STL-10 photographs
resized to $512\times512$. The retained scores are $0.997$, $0.853$, $0.733$ and $0.430$ for
$f=2,4,8,16$. The best gain achievable over the seven strengths is $-0.189$, $-0.061$, $-0.051$
and $+0.115$ respectively: at the three milder degradations no tested strength beats leaving the
image alone, and transport only begins to pay between $f=8$ and $f=16$, that is once the source's
alignment with its target falls below about $0.64$.

One scope note that belongs with these numbers: the metric is a single-target cosine, which is
the family the paper argues elsewhere should not by itself decide questions of this kind. The
experiment is reported as a targeted test on one checkpoint, one degradation family and $64$
images, not as a distributional validation.
\end{item12}

% =================================================================
\begin{item12}{3 and 4. The marginal-separation parameter, the regression, and the asymptotic
slope}
\ask{Specify the marginal-similarity parameter corresponding to the ratio sequence, provide the
exact regression giving slope $-0.756$, and insert the derivation of the asymptotic slope $-1$.
Send \texttt{ratio.tex}.}

\resp \texttt{ratio.tex} is attached. The sweep interpolates the second marginal toward the
first, $\Sigma_1(\delta)=\Sigma_0+\delta(\Sigma_1-\Sigma_0)$, so $\delta$ is the fraction of the
original separation retained; $\delta=1$ is the ensemble as drawn and $\delta\to0$ is the
homogeneous limit. Ensemble: dimension $16$, $64$ systems, relative violation $0.05$, median
response per unit violation, float64, second-order integration at $1{,}500$ steps.

\begin{center}
\small
\begin{tabular}{@{}lcccccc@{}}
\toprule
$\delta$ & $0.05$ & $0.1$ & $0.2$ & $0.4$ & $0.7$ & $1.0$\\
$\|\Sigma_1-\Sigma_0\|_F/\|\Sigma_0\|_F$ & $0.043$ & $0.087$ & $0.173$ & $0.347$ & $0.607$ &
$0.867$\\
skew/span ratio & $76.8$ & $42.0$ & $24.2$ & $14.6$ & $10.1$ & $7.8$\\
\bottomrule
\end{tabular}
\end{center}

The regression is ordinary least squares of $\log_{10}(\text{ratio})$ on $\log_{10}\delta$ over
these six points, giving slope $-0.756$ with rms residual $0.0166$.

\emph{The asymptotic slope.} Two facts give $-1$. The span response carries the factor
$\sup_t|(\log v_i/v_j)'|$ from the suppression bound in \texttt{ratio.tex}; that factor is
$O(\delta)$ when the marginals differ by $\delta$, so the span response vanishes linearly. The
skew response has no such factor and tends to the finite non-zero limit given in closed form by
the homogeneous-marginal proposition, so it is $O(1)$. The ratio is therefore $O(1/\delta)$ and
the log--log slope tends to $-1$.

The measured $-0.756$ is milder than $-1$ because the tested window is not yet asymptotic. The
ratio slope is identically the skew slope minus the span slope, an identity that holds in the
data to $10^{-15}$, so the discrepancy can be attributed exactly. Extending the sweep a decade
below the window above:

\begin{center}
\small
\begin{tabular}{@{}lccc@{}}
\toprule
window & span exponent & skew exponent & ratio slope\\
\midrule
$\delta\in[0.05,1]$     & $0.573$ & $-0.159$ & $-0.732$\\
$\delta\in[0.005,0.05]$ & $0.933$ & $-0.019$ & $-0.952$\\
\bottomrule
\end{tabular}
\end{center}

The span response is what is sublinear over the reported window, at $\delta^{0.57}$; the mild
decline of the skew constant contributes $-0.16$, which pushes the slope \emph{toward} $-1$ and
so cannot explain the shortfall. A decade lower the span exponent has moved to $0.93$ and the
ratio slope to $-0.95$, which is the predicted convergence. Slopes are ensemble-dependent at
about the $0.02$ level: an independently drawn pair of marginals gives $-0.756$ where the run
above gives $-0.732$.

The value $6.585$ should therefore be reported as the sampled ensemble's ratio at its own
separation, not as a constant. It grows without bound as the marginals merge.
\end{item12}

% =================================================================
\begin{item12}{5. The full first-order perturbation derivation}
\ask{Insert the full derivation, including the $t(1-t)$ suppression, the interpolant-variance
separation factor, the homogeneous-marginal first-order limit, and the closed-form arctangent
skew response.}

\resp All four are in \texttt{ratio.tex}, attached. In outline.

The response operator is $\mathcal D(E)=M_0\int_0^1\Phi^{-1}K_E\Phi\,dt$ with $K_E$ as in the
first-order proposition. The single identity that makes everything computable is that for any
symmetric expansion point, $G_0=\tfrac12 V_0'$.

\emph{Homogeneous marginals.} With $\Sigma_0=\Sigma_1=\Sigma$ and $C_0=\kappa\Sigma$, the clean
generator is a scalar multiple of the identity, so $M_0=I$ and the conjugation cancels. Every
symmetric direction then gives an integrand that is odd under $t\mapsto1-t$ against an even
weight, so $\mathcal D(S_E)=0$ \emph{exactly}, in or out of the span. Every skew direction gives
$\mathcal D(N_E)=-I_\kappa N_E\Sigma^{-1}$ with
$I_\kappa=\int_0^1 dt/\omega_\kappa(t)$, $\omega_\kappa=1-2(1-\kappa)t(1-t)$, which integrates to
$\tfrac{4}{\sqrt{\mu(4-\mu)}}\arctan\sqrt{\mu/(4-\mu)}\ge1$ for $\mu=2(1-\kappa)$. So at equal
marginals the ratio is infinite.

\emph{The two suppression factors.} In the commuting model the response acts pairwise, with the
skew part weighted by $\mathcal B_{ij}=\int_0^1 dt/\sqrt{v_iv_j}$, a strictly positive kernel
admitting no cancellation, and the symmetric part weighted by
$\mathcal A_{ij}=-\tfrac12\int_0^1 t(1-t)\,(\log v_i/v_j)'\,dt/\sqrt{v_iv_j}$. The symmetric
weight is the skew weight multiplied by the window $t(1-t)\le\tfrac14$, which vanishes at both
endpoints, and by the rate at which the two interpolant variances separate, which vanishes
identically when the marginals are homogeneous. Hence
$\mathcal B_{ij}/|\mathcal A_{ij}|\ge 8/\sup_t|(\log v_i/v_j)'|$.

A skew violation perturbs the drift directly; a symmetric violation reaches the endpoint map only
through the variance channel, which the parity of the clean flow cancels except to the extent
that coordinates traverse different variance profiles. That is the whole mechanism.

Verification: six numerical checks in \texttt{notebooks/ratio\_proof.py}, all holding. The
pairwise formulas match a finite-differenced endpoint map to $1.3\times10^{-6}$ over $33$
perturbations; the two forms of $\mathcal A_{ij}$ agree to $2.8\times10^{-17}$; on pairs with
proportional variance profiles the symmetric response measures below $9.1\times10^{-9}$; and with
homogeneous marginals the symmetric response is $3.5\times10^{-10}$ against a skew response of
$1.54$ matching the closed form to $3.2\times10^{-9}$.
\end{item12}

% =================================================================
\begin{item12}{6. Whether $\Delta$ can have a root when $A\ge\cos\jmath$
\;---\; \textcolor{red}{this is a correction}}
\ask{The endpoint signs give sufficiency only. Can you prove $\Delta$ has no root in $(0,1)$ when
$A\ge\cos\jmath$?}

\resp No, because it is false. $\Delta$ always has a root. The stated equivalence should be
withdrawn and replaced; the corrected statement is stronger, and it recovers three grid cells the
sweep had been discarding. Full proof in \texttt{threshold.tex}, attached; every step verified
numerically over $120$ cells with no exceptions.

Substituting $\beta=\sin\varphi$ and writing $\psi(\varphi)=\arccos(A\cos\varphi)$ turns both
reference levels into cosines of angle differences:
\[
\Delta/A=\cos\!\big(\psi(\varphi)-\varphi\big)-\cos(\varphi-\jmath),
\]
an identity holding to $7\times10^{-15}$. At $\varphi=\jmath$ the comparison term is $\cos 0=1$,
its maximum, while $\psi(\jmath)>\jmath$ strictly for any $A<1$. Hence
\[
\Delta(\sin\jmath)<0\quad\text{for every }A\in(0,1),\ \jmath\in(0,\pi/2),
\]
irrespective of $A$. Since $\Delta(1)=A(1-\sin\jmath)>0$, a transition above $\sin\jmath$ exists
unconditionally.

Uniqueness is what $A<\cos\jmath$ governs. Above $\sin\jmath$ the sign of $\Delta$ is the sign of
$h(\varphi)=2\varphi-\jmath-\psi(\varphi)$, and since $\psi'=A\sin\varphi/\sin\psi\le A$ we have
$h'\ge 2-A>1$, so $h$ is strictly increasing and crosses zero exactly once. Below $\sin\jmath$
the condition reduces to $A\cos\varphi>\cos\jmath$, strictly decreasing in $\varphi$, so a second
root appears exactly when $A>\cos\jmath$. Therefore
\[
A<\cos\jmath\iff\Delta\text{ has a single sign change on }(0,1),
\]
and when $A>\cos\jmath$ the flow wins on $(0,r_1)$, loses on $(r_1,r_2)$ and wins again on
$(r_2,1)$, with $r_1<\sin\jmath<r_2$.

The consequence for the table is that the three cells marked ``--'' are not empty. The sweep
already measured their upper roots, and those measurements are as accurate as the ones reported:

\begin{center}
\small
\begin{tabular}{@{}rrrrr@{}}
\toprule
$A$ & $\jmath$ & $\sin\jmath$ & analytic $r_2$ & measured\\
\midrule
$0.70$ & $0.80$ & $0.7174$ & $0.8347$ & $0.8286$\\
$0.90$ & $0.50$ & $0.4794$ & $0.5869$ & $0.5600$\\
$0.90$ & $0.80$ & $0.7174$ & $0.7708$ & $0.7586$\\
\bottomrule
\end{tabular}
\end{center}

The comparison therefore covers $30$ cells, not $27$, and the registered entry recording three
failures of the existence criterion should be withdrawn: the criterion was wrong, the measurement
was right.
\end{item12}

% =================================================================
\begin{item12}{7. What ``$4\times$ training'' means}
\ask{Specify precisely what $4\times$ training means, and confirm all remaining settings are
unchanged.}

\resp Optimisation steps: $3{,}000$ against $12{,}000$, exactly four times as many. Everything
else is held fixed --- batch $256$, AdamW at $3\times10^{-4}$ under maximal-update
parameterisation, gradient clipping at $1.0$, the same checkpoint-selection rule, and the same
evaluation protocol and seeds. The long run differs from the short one in the step-count argument
alone.
\end{item12}

% =================================================================
\begin{item12}{8. The capacity sweep at a second learning rate}
\ask{Confirm whether the capacity sweep has been repeated at a second learning rate.}

\resp It has been run at three: $\eta\in\{10^{-4},3\times10^{-4},10^{-3}\}$ across four widths
with five seeds each, $120$ trainings under the same protocol. The gap $\M_1-\M_2$:

\begin{center}
\small
\begin{tabular}{@{}rrrrr@{}}
\toprule
$\eta$ & width $256$ & width $512$ & width $1024$ & width $2048$\\
\midrule
$10^{-4}$        & $0.0419$ & $0.0524$ & $0.0558$ & $0.0612$\\
$3\times10^{-4}$ & $0.0835$ & $0.0938$ & $0.0873$ & $0.0767$\\
$10^{-3}$        & $0.0892$ & $0.1158$ & $0.1020$ & $0.1018$\\
\bottomrule
\end{tabular}
\end{center}

The sign is the same in all twelve cells, at a minimum of $16.2$ standard errors. The mean gap
grows monotonically with the learning rate at every width ($0.0528$, $0.0853$, $0.1022$; Spearman
$0.90$). Raising the step size from $10^{-4}$ to $10^{-3}$ costs the marginal model $1$--$7\%$ of
its terminal fidelity and the source-conditioned model $19$--$27\%$: the source-conditioned arm
is four to sixteen times more sensitive to step size at matched architecture and objective, which
is what an amplification near the initial time predicts and a generic optimisation failure does
not.

Two things to report rather than round away. At no learning rate does the widest model close the
gap at the two-standard-error level, but at $3\times10^{-4}$ the gap does narrow with width,
$0.0835\to0.0767$, reaching $1.9$ standard errors; the claim should be ``does not close'', not
``widens''. And the diversity collapse is universal across all twelve cells and worsens with
capacity: the source-conditioned pairwise cosine runs $0.344\to0.506$ from width $256$ to $2048$
at $\eta=10^{-4}$, against a marginal-model mean of $0.225$ and a true-posterior reference of
$0.050$.
\end{item12}

% =================================================================
\begin{item12}{9. The output-variety metric, the step sweep, and the integration scheme}
\ask{Provide the exact definition of the output-variety metric, its values over the complete
integration-step sweep, the tested step counts, and the numerical integration scheme. Confirm the
one-shot regression uses the same evaluation set and target-alignment definition.}

\resp Variety is measured as the mean pairwise cosine among the six samples generated for one
condition, averaged over conditions. The direction matters: this number \emph{decreases} as
variety increases. The solver is Heun's method, second order, intrinsic to the sphere, on a
uniform time grid, with $32$ steps by default. Step counts tested: $K\in\{1,2,4,8,16,32,64\}$.

\begin{center}
\small
\begin{tabular}{@{}rrrr@{}}
\toprule
$K$ & target cosine & pairwise cosine & polar $W_1$\\
\midrule
$1$  & $0.3541$ & $0.5299$ & $0.4623$\\
$2$  & $0.3524$ & $0.3897$ & $0.3286$\\
$4$  & $0.3389$ & $0.2593$ & $0.1632$\\
$8$  & $0.3254$ & $0.2030$ & $0.0611$\\
$16$ & $0.3197$ & $0.1858$ & $0.0360$\\
$32$ & $0.3175$ & $0.1838$ & $0.0327$\\
$64$ & $0.3172$ & $0.1829$ & $0.0321$\\
\bottomrule
\end{tabular}
\end{center}

A single integration step attains the highest target cosine and the worst distributional fit, by
a factor of fourteen in $W_1$; sixty-four steps invert both. Selecting the integration budget by
single-target cosine alone therefore selects the coarsest possible rollout.

Confirmed: the one-shot regression baseline is scored by the identical routine on the identical
task, $768$ conditions with $6$ samples each, and the identical target-alignment definition.
\end{item12}

% =================================================================
\begin{item12}{10. The classifier-free-guidance setup}
\ask{Provide the exact formulation, all tested guidance weights, the definition and tested values
of $\lambda$, the underlying terminal-cosine values, and the unguided reference used for the
reported $+0.157$ and $-0.505$.}

\resp The guided velocity is
$v_w=v_{\mathrm{uncond}}+w\,(v_{\mathrm{cond}}-v_{\mathrm{uncond}})$, tangent-projected to the
sphere. The unconditional branch is the same network with its condition replaced by a learned
null token, trained with $10\%$ condition dropout. Weights tested: $w\in\{1,2,3\}$.

\emph{The unguided reference is $w=1$}, where the formula returns the conditional field
unchanged. Every reported change is measured against it, and the $w=1$ column below is flat at
zero to within $0.0014$, which is the internal consistency check.

$\lambda$ interpolates the initial state from noise toward the informative source,
$z_0\leftarrow\mathrm{slerp}(u,z_0,\lambda)$ with $u$ uniform on the sphere: $\lambda=0$ starts
from noise, $\lambda=1$ from the fully informative source. Tested at
$\lambda\in\{0,0.25,0.5,0.75,1\}$, two seeds per cell. Change in terminal target cosine relative
to $w=1$:

\begin{center}
\small
\begin{tabular}{@{}rrrr@{}}
\toprule
$\lambda$ & $w=1$ & $w=2$ & $w=3$\\
\midrule
$0.00$ & $+0.0001\pm0.0012$ & $+0.0424\pm0.0002$ & $+0.0770\pm0.0004$\\
$0.25$ & $-0.0000\pm0.0012$ & $-0.0325\pm0.0003$ & $-0.0673\pm0.0007$\\
$0.50$ & $+0.0001\pm0.0014$ & $-0.1271\pm0.0003$ & $-0.2244\pm0.0001$\\
$0.75$ & $+0.0002\pm0.0014$ & $-0.2430\pm0.0138$ & $-0.3346\pm0.0065$\\
$1.00$ & $+0.0007\pm0.0011$ & $-0.3721\pm0.0173$ & $-0.4905\pm0.0039$\\
\bottomrule
\end{tabular}
\end{center}

On the reported values: $+0.157$ and $-0.505$ come from an earlier run. The current values, with
seed error bars, are $+0.0770\pm0.0004$ and $-0.4905\pm0.0039$, and the draft should use these.
The sign changes at $\lambda_0=0.13$, not near $\lambda=1$: guidance calibrated for a noise start
begins to hurt as soon as the initialisation carries even slight target information, which is a
sharper statement than the one currently in the text.
\end{item12}

% =================================================================
\begin{item12}{11. The OpenCLIP protocol, and the factorization discrepancy $0.62$
\;---\; \textcolor{red}{contains a correction}}
\ask{Provide the architecture and checkpoint, the CIFAR-10 split and image count, preprocessing
and blur construction, embedding normalisation, the target-alignment metric, the class prototype,
and the precise definition and computation of the factorization discrepancy $0.62$.}

\resp OpenCLIP ViT-B-32, checkpoint \texttt{laion2b\_s34b\_b79k}. CIFAR-10 \emph{train} split,
$6{,}000$ images drawn with a fixed generator seed. Preprocessing is the transform shipped with
that checkpoint. Image embeddings are $L_2$-normalised. The source is the embedding of a
Gaussian-blurred image, kernel size $9$, with $\sigma$ drawn uniformly from $[0.8,3.0]$ per call.
The condition is the $L_2$-normalised mean of the clean embeddings of the image's own class.
Split $80/20$; the flow has width $512$ and three blocks, trained $3{,}000$ steps at
$\eta=10^{-4}$ with batch $256$, sampled with $32$ Heun steps. Target alignment is the mean
cosine between the output embedding and the clean embedding of the same image, on the held-out
$20\%$.

Measured levels: random pair $0.4815$, blurred source $0.5517$, marginal model $0.7586$,
source-conditioned $0.7507$, class prototype $0.7930$. The estimated parameters are
$\hat A=0.7923$, $\hat\beta=0.6010$, $\hat\jmath=0.2817$, giving a predicted threshold of
$0.5291$, below $\hat\beta$; the prediction that the flow beats the source was therefore made in
advance and holds, at $89\%$ of the range between the random-pair floor and the prototype. The
model ordering also transfers.

\emph{On the $0.62$.} No quantity of that name is computed anywhere in the pipeline. The stage
was re-run to check, and its output contains no column equal to $0.62$; the closest value is
$\hat\beta=0.6010$, which already appears in the paper under its own name and role. The sentence
should be removed unless it can be traced to a specific quantity.

A related item worth recording: the label probe in this section is now trained and scored on
disjoint $80/20$ splits and reaches $0.058$ held-out accuracy against a chance level of $0.10$,
so it does not extract label information from the private residual. An earlier version was
trained and scored on the same $6{,}000$ embeddings and was withdrawn as instrumentation; the
current one restores the result properly.
\end{item12}

% =================================================================
\begin{item12}{12. The completed super-resolution dimensionality experiment}
\ask{Insert the completed experiment, including the three resolutions, the theoretical
dimensional bounds, the measured source and target dimensions, and sufficient detail of the
estimation procedure to reproduce it.}

\resp STL-10, $5{,}000$ training and $1{,}024$ test images, resized and centre-cropped to each
resolution and normalised to $[-1,1]$. The condition $c$ is the class label; the source is
$z_0=\mathsf{UD}z_1$, average pooling by $f=4$ followed by nearest-neighbour upsampling, so that
conditioned on $c$ the source is a deterministic function of the target and the endpoints are
maximally dependent. $\mathsf P=\mathsf{UD}$ is an exact projection: the identity arm carries
relative null-space energy $7.5\times10^{-15}$.

\emph{Estimation procedure.} Rank is read from the eigenvalues of the Gram matrix of the
mean-centred sample, computed in float64, with a relative cutoff of $10^{-10}$ against the
largest eigenvalue. The sample size is chosen as $\min(5000,\,4r+256,\,3072)$ where $r$ is the
predicted bound, so that the number of samples strictly exceeds the bound at every resolution and
no figure is sample-limited from below. That gives $1{,}024$, $1{,}984$ and $3{,}072$ samples at
the three resolutions respectively.

\begin{center}
\small
\begin{tabular}{@{}rrrrrr@{}}
\toprule
resolution & ambient & bound $r=3(\mathrm{res}/f)^2$ & measured rank $z_0$ &
measured rank $z_1$ & samples\\
\midrule
$32$ & $3072$  & $192$  & $192$  & $1023$ & $1024$\\
$48$ & $6912$  & $432$  & $432$  & $1983$ & $1984$\\
$96$ & $27648$ & $1728$ & $1728$ & $3071$ & $3072$\\
\bottomrule
\end{tabular}
\end{center}

The measured source rank equals the predicted bound exactly at all three resolutions. The target
ranks are one below the sample count, that is they are limited by the number of samples and not
by the data; what the table establishes is the strict inequality
$\dim\mathrm{supp}(z_0)\le 192 < 1023 \le \dim\mathrm{supp}(z_1)$ and its analogues, which is
what the obstruction argument requires. It does not establish the true target dimension, and the
claim should be stated in that form.

The same measurement was repeated at fixed resolution $48$ across six downsampling factors, and
is exact at every one:

\begin{center}
\small
\begin{tabular}{@{}rrrrrrr@{}}
\toprule
factor & $2$ & $3$ & $4$ & $6$ & $8$ & $12$\\
\midrule
predicted bound & $1728$ & $768$ & $432$ & $192$ & $108$ & $48$\\
measured rank   & $1728$ & $768$ & $432$ & $192$ & $108$ & $48$\\
\bottomrule
\end{tabular}
\end{center}

Nine exact agreements in total. One qualification to state plainly: the bound is an upper bound
by construction, since $z_0$ lies in the range of a fixed rank-deficient operator, so what is
being confirmed is that natural images saturate it rather than that a non-trivial prediction was
risked.

At those six factors the trained flow beats returning the source at every level, with the gain
falling as the source degrades --- $3.81$, $2.26$, $1.97$, $1.29$, $1.02$, $0.59$ PSNR. This is
the opposite direction from the pretrained image-to-image experiment of item~2, and the contrast
is the useful part: a flow trained for its exact degradation beats doing nothing everywhere we
tested, while a pretrained flow applied off the shelf is worse than doing nothing at $2\times$,
$4\times$ and $8\times$. The do-nothing regime lives in the mismatch, not in the task. At the
harshest factor the least-squares regression arm does fall below the baseline, at $-0.32$ PSNR.
\end{item12}

\end{document}
